\subsection{固定正定度量的特例}
\label{sec:gppa-transfer}

固定对称正定矩阵 $M$，定义
\[
 \begin{aligned}
 \langle h,k\rangle_M:=\langle Mh,k\rangle,\qquad
 &\|h\|_M:=\sqrt{\langle Mh,h\rangle},\\
 d_M(x,S):=\inf_{p\in S}\|x-p\|_M,\qquad
 &r_{F,M^{-1}}(y):=\inf_{f\in F(y)}\|f\|_{M^{-1}}.
 \end{aligned}
\]
考虑预条件近端关系
\[
 J^M_{\lambda F}(x):=(M+\lambda F)^{-1}(Mx).
\]
其每个输出 $y$ 满足 $M(x-y)\in\lambda F(y)$。
正定性使 $\|\cdot\|_M$ 成为与欧氏范数等价的范数。
以下推论因此保留前述收敛证明，但所有条件均须在这个范数下检验。

\begin{corollary}[固定正定核的局部收敛]
\label{cor:gppa-spd}
沿用推论~\ref{thm:selection-RLEB} 的对象与参数，并令
$U_{M,R}=\{x\in U:d_M(x,S)\leq R\}$。
将该定理中的近端关系换成 $J^M_{\lambda F}$，
内积换成 $\langle\cdot,\cdot\rangle_M$，
范数、距离和最近零点分别换成 $M$ 范数、$d_M$ 和 $M$ 最近零点，
真实残差换成 $r_{F,M^{-1}}$，误差界函数记为 $\psi_M$。
特别地，(B2) 等价于
$\|2y-x-p\|_M\leq Ld_M(x,S)^\gamma$，其中 $p$ 是输入的 $M$ 最近零点。
若允许关系 $\mathcal J_M$ 在这些替换下满足 (B1)--(B4) 及严格初始长度预算，
则该定理的全部结论在 $M$ 范数下成立。
特别地，取
\[
 \mathcal L_M(r)=\frac12\left(
 \frac{r}{1-\kappa}+\frac{Lr^\gamma}{1-\kappa^\gamma}\right),
 \qquad d_0=d_M(x^0,S),
\]
则所有允许轨道具有有限长度并收敛到 $S\cap U$ 中的一点，且
\[
 d_M(x^k,S)\leq\kappa^k d_0,\qquad
 \|x^\infty-x^k\|_M\leq\mathcal L_M(d_M(x^k,S)).
\]
若 $m=\lambda_{\min}(M)$，欧氏总长度至多
$\mathcal L_M(d_0)/\sqrt m$，欧氏点尾至多
$\mathcal L_M(\kappa^k d_0)/\sqrt m$。
当 $\mathcal J_M=J^M_{\lambda F}$ 时，结论控制全部完整近端选择。
\end{corollary}

\begin{proof}
置 $Q=M^{1/2}$，$z=Qx$，$Z=QS$，$W=QU$，并定义完整关系
\[
 G(z):=Q^{-1}F(Q^{-1}z).
\]
由于 $Q$ 是线性同胚，$Z$ 非空闭、$W$ 开。
对任意输入 $x$ 和任意输出 $y$，
\[
 M(x-y)\in\lambda F(y)
 \quad\Longleftrightarrow\quad
 Qx-Qy\in\lambda G(Qy).
\]
因此两个完整近端关系逐输入、逐输出对应。
同时，由定义直接有
\[
 d(Qx,Z)=d_M(x,S),\qquad
 r_G(Qy)=r_{F,M^{-1}}(y),\qquad
 d(Qx,\mathbb R^n\setminus W)
 =d_M(x,\mathbb R^n\setminus U).
\]
最近零锚及反射估计也在 $Q$ 下逐项对应。
故推论的假设正是推论~\ref{thm:selection-RLEB}
对 $G,Z,W$ 及允许关系
$\widetilde{\mathcal J}(z)=Q\mathcal J_M(Q^{-1}z)$ 的假设。
应用该定理再用 $Q^{-1}$ 回代，得到全部 $M$ 范数结论。
最后，$\|h\|\leq\|h\|_M/\sqrt m$ 给出欧氏长度与点尾估计。
\end{proof}

更换核时，覆盖、反射估计及兼容条件必须重新检验。
这一点在关系 $F(s,t)=(s,-\sqrt{|s|})$ 上已经可见。
取 $\lambda=1$ 和 $M=\operatorname{diag}(a,b)$，$a,b>0$，
记 $q=a/(a+1)$。完整近端关系逐输入单值，且
\[
 J^M_F(s,t)=\left(qs,t+\frac{\sqrt{q|s|}}b\right).
\]
因此每条轨道的法向距离按 $q$ 收缩，切向增量可和；
其极限是
\[
 \left(0,t^0+\frac{\sqrt{q|s^0|}}{b(1-\sqrt q)}\right),
\]
欧氏总长度至多 $|s^0|+\sqrt{q|s^0|}/[b(1-\sqrt q)]$。
这些结论由显式公式直接得到，不要求先调用一般推论。

正定性本身却不足以保证同样的结论。仍取这个 $F$，令
\[
 M=\begin{pmatrix}1&1/2\\1/2&1\end{pmatrix}.
\]
该矩阵正定，而输入 $0$ 的完整近端关系同时包含
$0$ 和 $(-4/49,16/49)$。
后一输出不属于零集，故上述零锚反射条件在输入 $0$ 已经失败。
这一例子说明：成功的正对角核计算不能直接推广到所有正定核。
